\section{A growing set of coefficient primes for the actual cubic characters}
\label{sec:mixed-cubic-split-prime-strict}

Put $K=\mathbb Q(\omega)$, $t=\sqrt[3]{269}$ and $L=K(t)$.
Choose either faithful character
$\chi:\operatorname{Gal}(L/K)\to\mu_3$; both powers $\chi^j$,
$j=1,2$, will be treated.  Define the explicit prime set
\[
 \mathscr P=\left\{\ell\text{ prime}:\ \ell>3,\quad
 \ell\ne7,269,\quad \ell\equiv1\pmod3,\quad
       269^{(\ell-1)/3}\not\equiv1\pmod\ell\right\}.
 \tag{MCP.1}
\]
The congruence tests in this definition are conditions on rational
integers.  In particular they do not add a finite character to the
actual host or require roots of unity of order $\ell$ in its
coefficient order.

For $\ell\in\mathscr P$, let $O$ be the original ring of integers
of any finite extension $E/\mathbb Q_\ell$ containing $\mu_3$,
with uniformizer $\varpi$ and residue field $k$.  Let
$\kappa:G_K\to1+\mathfrak m_O$ be continuous and unramified
away $\ell$.  For $e\geq1$ divisible by $\ell-1$, put
\[
        D=(E/O)(\chi^j\kappa^{-1}c_\ell^{-e}),\qquad j=1,2.
 \tag{MCP.2}
\]
Here $c_\ell$ is the usual cyclotomic character.  Define the full
strict group inside $H^1(G_K,D)$ by zero restriction to $H^1(K_v,D)$ for every
$v\mid\ell$, and to $H^1(I_v,D)$ at every other finite place.
Use the identical conditions for the finite module $D[\varpi]$.

\begin{theorem}[AC-G203: original-coefficient strict zero for a growing mixed family]
\label{AC-G203}
For all the choices in (MCP.1)--(MCP.2), on the original $O$,
\[
                \mathcal S_\ell(K,D[\varpi])=0,
                \qquad \mathcal S_\ell(K,D)=0.
 \tag{MCP.3}
\]
The same conclusion holds for the actual groups (FC.2), with
their unchanged full local conditions, at the two original
finite labels $\eta=\chi$ and $\eta=\chi^2$ and these weights.
All rational two-power, conductor-seven, and auxiliary $q,d$
character factors of these labels are trivial.

For every integer $T\geq2$, the exact set of such original triples is
\[
 \begin{split}
 \mathcal M(T)=\{(\ell,\chi^j,e):\ &\ell\in\mathscr P,
     \ j=1,2,\ 1\leq e<T,\ \ell-1\mid e\},\\
 |\mathcal M(T)|={}&2\sum_{\substack{\ell\in\mathscr P\\\ell\leq T}}
                  \left\lfloor\frac{T-1}{\ell-1}\right\rfloor.
 \end{split}
 \tag{MCP.4}
\]
The coefficient primes $13$ and $19$ belong to $\mathscr P$.
Their disjoint contribution to this count is
\[
 2\left\lfloor\frac{T-1}{12}\right\rfloor
       +2\left\lfloor\frac{T-1}{18}\right\rfloor.
 \tag{MCP.5}
\]
More generally, for any retained set $J$ of actual triples,
the nonnegative remainder of (CP.3) satisfies the exact identity
\[
                       \mathcal R_J
                          =\mathcal R_{J\setminus\mathcal M(T)}.
 \tag{MCP.6}
\]
No derived determinant index is asserted to vanish.  Equation
(MCP.6) is an equality of the actual strict-group sum before
the separate full-union index comparison and window transport.
\end{theorem}

\begin{proof}
\proofstep{1. The fixed field and the completed class-quotient certificate.}
The polynomial $X^3-269$ is Eisenstein at $269$, so
$[\mathbb Q(t):\mathbb Q]=3$.  Its intersection with the quadratic
field $K$ is $\mathbb Q$, by the degrees.  Consequently
$[L:\mathbb Q]=6$ and $\operatorname{Gal}(L/K)=C_3$.
For $z=t+\omega$ the exact norm polynomial is
\[
 \begin{split}
 P(X)&=\bigl((X-\omega)^3-269\bigr)
          \bigl((X-\omega^2)^3-269\bigr)\\
 &=X^6+3X^5+6X^4-531X^3-801X^2+810X+72900.
 \end{split}
 \tag{MCP.7}
\]
Its six conjugate roots are distinct.  An equality between two
of $t\omega^i+\omega^a$, with $i\in\{0,1,2\}$ and $a\in\{1,2\}$,
would either immediately identify both indices or express $t$
as an element of $K$.  The latter is impossible by its degree.
Thus $z$ generates exactly $L$, proving the identity of the field
used in the computation, not just its degree or signature.

The completed root-owned PARI/GP input constructs
$B=\operatorname{bnfinit}(P,1)$ and then runs
$\operatorname{bnfcertify}(B,1)$.  Its saved exact output is
\[
 \begin{array}{c|c}
 B.\mathrm{sign}&[0,3]\\
 B.\mathrm{disc}&-141375086667\\
 B.\mathrm{no}&49\\
 B.\mathrm{cyc}&[7,7]\\
 \operatorname{bnfcertify}(B,1)&1.
 \end{array}
 \tag{MCP.8}
\]
The input, complete output and saved $B$ are retained verbatim.
The final output also records completion after the successful
certificate.  The documented flag-one conclusion is precisely
that the actual class group is a quotient of the computed group.
It follows that
\[
        \operatorname{Cl}(L)\text{ is a quotient of }C_7^2,
        \qquad \operatorname{Cl}(L)/\ell=0\quad(\ell\ne7).
 \tag{MCP.9}
\]
Only this conclusion is used.  The computation alone does not
assert equality with $C_7^2$, a nonzero seven-primary class group,
or that the displayed units form a full basis.  The complete
PARI manual clause and flag-one interface are retained as the
software-certificate input; no new proof of the whole algorithm
is claimed.  The author inspected this completed root execution
and did not rerun it.

\proofstep{2. Residual coefficients and both split local Frobenius elements.}
The reduction of $c_\ell$ has values in
$\mathbb F_\ell^\times$ and is killed by its $(\ell-1)$st power.
The specified $\kappa$ reduces to one.  Reduction on $\mu_3$
is injective because $\ell>3$.  Hence
\[
                         D[\varpi]=k(\overline\chi^{\,j}),
 \tag{MCP.10}
\]
with a character of exact order three and splitting field $L/K$.
This uses the actual residue field $k$ and does not extend or
normalize the original coefficient ring.

Since $\ell\equiv1\pmod3$, the polynomial $X^2+X+1$ has
two distinct roots modulo $\ell$.  Thus there are two places
$v$ of $K$ over $\ell$, and each completion is $\mathbb Q_\ell$.
At each such place, $X^3-269$ has separable reduction because
$\ell\ne3,269$.  The multiplicative group of the residue field
is cyclic, and the last condition in (MCP.1) says exactly that
$269$ is not a cube in that field.  A cubic polynomial is
reducible over a field only if it has a root.  Therefore this
polynomial is irreducible modulo $\ell$ at both places.
The local extension $L_w/K_v$ is unramified of degree three,
and arithmetic Frobenius generates its cyclic Galois group.
In particular
\[
       \overline\chi^{\,j}(\operatorname{Frob}_v)\ne1
                       \qquad(j=1,2;\ v\mid\ell).
 \tag{MCP.11}
\]
Changing the selected cubic character or the Frobenius convention
only inverts a nontrivial cube root; it does not change this fact.

\proofstep{3. Vanishing of the full strict residual group.}
Write $A=k(\overline\chi^{\,j})$.
Inflation--restriction for $G_L\subset G_K$ makes
\[
                       H^1(G_K,A)\longrightarrow H^1(G_L,k)
 \tag{MCP.12}
\]
injective: its kernel is $H^1(C_3,A)=0$, since averaging by
three is valid in characteristic $\ell>3$.
The coefficient on $G_L$ is trivial, so every restricted
continuous cocycle is an additive homomorphism into $k$.
For a strict residual class it kills all finite inertia groups
of $L$.  At places above $\ell$ it even kills the full local
group, since the original local class over $K$ was zero.
At other places its inertia restriction remains zero under
restriction to $L$.  There are no real places.
Unramified abelian class field theory therefore factors this
homomorphism through $\operatorname{Cl}(L)$.  Equivalently there
is an injection into the appropriate equivariant part of
\[
                    \operatorname{Hom}(\operatorname{Cl}(L),k)=0,
 \tag{MCP.13}
\]
where the zero follows from (MCP.9).  This proves the first
assertion in (MCP.3).  No character-orientation convention in
Artin equivariance affects this zero, since the whole
$\ell$-primary class group is absent.

\proofstep{4. The original local invariant quotients all vanish.}
For either $v\mid\ell$, choose a Frobenius lift in $G_{K_v}$.
Its scalar on $D$ reduces to the nontrivial value in (MCP.11),
even when $\kappa$ is ramified at $v$.  Its scalar minus one
is consequently an $O$-unit.  Multiplication by this unit is
an automorphism of $E/O$, proving
\[
                     D^{G_{K_v}}=0\quad(v\mid\ell),
                     \qquad D^{G_K}=0.
 \tag{MCP.14}
\]
This checks both split local conditions separately.

Let $v\nmid\ell$ be any other finite place.  Both $c_\ell$
and $\kappa$ are unramified there.  The image of
$\chi^j(I_v)$ is either trivial or has order three.
In the nontrivial case some inertia scalar minus one is an
$O$-unit, so $D^{I_v}=0$.  In the trivial case inertia acts
trivially on all of $D$, so $D^{I_v}=D$ and multiplication by
$\varpi$ is surjective.  This includes the possibly ramified
cubic places over three and $269$, without enlarging or
weakening their local conditions.  We have proved
\[
 \frac{D^{G_{K_v}}}{\varpi D^{G_{K_v}}}=0\ (v\mid\ell),
 \qquad
 \frac{D^{I_v}}{\varpi D^{I_v}}=0\ (v\nmid\ell).
 \tag{MCP.15}
\]

\proofstep{5. Lift every original-uniformizer strict class.}
Use the exact sequence of the original $O$-modules
\[
       0\longrightarrow D[\varpi]\longrightarrow D
                      \xrightarrow{\varpi}D\longrightarrow0.
 \tag{MCP.16}
\]
For every global or local group $G$ under discussion, its
cohomology sequence contains
\[
  D^G/\varpi D^G\longrightarrow H^1(G,D[\varpi])
               \longrightarrow H^1(G,D)[\varpi]\longrightarrow0.
 \tag{MCP.17}
\]
Take a class $x\in\mathcal S_\ell(K,D)[\varpi]$.
Global surjectivity in (MCP.17) gives a residual lift $\bar x$.
At each prescribed local group its image in $H^1(G,D)$ is zero.
The local exact sequence and (MCP.15) imply that its local
class in $H^1(G,D[\varpi])$ is already zero.  Thus $\bar x$
is a full strict residual class, which vanishes by Step 3.
It follows that $\mathcal S_\ell(K,D)[\varpi]=0$.

Every continuous cocycle from the compact profinite group $G_K$
to the discrete module $D$ has finite image, hence is killed
by some power of $\varpi$.  Any nonzero cohomology class in
the strict group would therefore give a nonzero class killed
by $\varpi$, after multiplication by a suitable power.
The preceding vanishing excludes this and proves the second
assertion in (MCP.3).  No unramified-outside-$\ell$ action of
the full coefficient over $L$ and no local generation theorem
for its maximal pro-$\ell$ extension is required.

\proofstep{6. Actual host, growing primes, and retained sums.}
The product decomposition in (FC.1) includes the factor
$\operatorname{Gal}(L/K)=C_3$.  Its two faithful characters with
every other factor trivial are exactly the two labels used here.
Their values lie in the prescribed completion $O_\ell$ of
$\mathbb Z[\zeta_{3\cdot2^h}]$ for every actual page $h\geq1$.
There is a canonical coefficient identification
\[
 (O_\ell\otimes_{\mathbb Z_\ell}\mathbb Q_\ell/\mathbb Z_\ell)
       (\chi^j)(-e)
                 =(E_\ell/O_\ell)(\chi^j c_\ell^{-e}).
 \tag{MCP.18}
\]
Apply the result with $\kappa=1$.  Inflation from the actual
$G_{K,\Sigma_\ell}$ to $G_K$ is injective, because the coefficient
is unramified outside that set.  Its strict classes are inertia-zero
outside that set and satisfy exactly the same specified local
conditions inside it.  They therefore inject into the zero group
in (MCP.3).  The original $O_\ell$, its uniformizer, and its
normalizing rank $b_\ell$ are all retained.

For fixed $\ell$ the allowed positive weights number
$\lfloor(T-1)/(\ell-1)\rfloor$.  If this is nonzero, then
$\ell\leq T$.  The two cubic labels are distinct; primes are
separate coordinates of the original triple set.  This proves
(MCP.4), without any density theorem or asymptotic for its
explicit prime criterion.  At $13$ one has $269\equiv9$ and
$9^4\equiv9\ne1$; at $19$ one has $269\equiv3$ and
$3^6\equiv7\ne1$.  These are the two examples in (MCP.5).

Every triple of $\mathcal M(T)$ has nontrivial cubic factor.
It therefore contributes to the first sum of (CP.3), where
the term just proved zero is the full original $K$-strict group,
not a rational parity or an individual derived index.
Deleting these zero summands from any $J$ proves (MCP.6).
There is no new coefficient-comparison error or independent
window in this identity.  An ensuing window theorem may apply
the existing one-union (FL.4) cost and the previously proved
critical comparison once, on that same prescribed window.
\end{proof}

\paragraph{Proof scope and earlier tests.}
The primary class-field input is Sharifi's Theorem 1.1.11, whose
complete statement is retained, and the complete 1.3.1--1.3.5 passage, including Artin
equivariance.  The exact original-uniformizer argument is written
out above; its prior instance is AC-G184.  PARI's completed
flag-one quotient certificate is the only class computation used.
The root proposed and computed the fixed mixed field and tested
local unit generation at $7$, $13$ and $19$; the present author
supplied the direct residual-lifting route and its growing-prime
scope.  Failure of that local-unit test at $13$ is not failure of
the theorem.  The successful $19$ unit test is an unused
alternative.  The class quotient alone settles nothing about
seven-primary nonvanishing, and $7$ is excluded throughout.
The two nontrivial finite cubic factors are genuine original
characters, so this result concerns an actual mixed family.
The audited finite-cover and Artin--Tate dictionary is unchanged;
no new motive, realization, or Lean execution is asserted.
The remaining all-prime arithmetic budget, $\mathcal R_{\rm cub}$
and H04.4 are not closed here.


\paragraph{Finite-cover and Artin--Tate dictionary.}
For every $m\geq1$, the original finite module $D[\varpi^m]$
defines a finite locally constant etale sheaf on
$\operatorname{Spec}\mathcal O_K[1/(3\cdot269\cdot\ell)]$:
its continuous finite action is unramified outside precisely the
allowed finite set.  The residual scalar action is trivialized by
the $C_3$ cover coming from $L/K$ on this open scheme.
The strict kernels used above retain restriction to each
$\ell$-adic completion and to every other prescribed inertia group;
these are the same original finite-coefficient cohomology conditions.
This is the finite-cover/Artin--Tate interpretation in the retained
SGA~1, Expose V, sections~4 and~7 audit used for the (FC.2)
dictionary.  No fresh archival reading, new motivic realization,
or replacement of the original integral sheaf is claimed.
